\begin{center}
\phantomsection
\textbf{\large{Tóm tắt}}
\end{center}

\addcontentsline{toc}{chapter}{Tóm tắt}

\begin{small}

Công bố thông tin ESG ngày càng đóng vai trò trung tâm trong báo cáo doanh nghiệp ngành ngân hàng, nhưng sự mở rộng nhanh chóng của hoạt động truyền thông ESG đã làm gia tăng lo ngại về hiện tượng ``ESG-washing''. Các phương pháp phát hiện hiện nay chủ yếu dựa vào xếp hạng ESG của bên thứ ba, vốn còn hạn chế về độ bao phủ và thiếu minh bạch.

Nghiên cứu này đề xuất một khung phương pháp neuro-symbolic dựa trên NLP nhằm đánh giá rủi ro ESG-washing thông qua khả năng kiểm chứng nội tại của các tuyên bố ESG trong báo cáo thường niên ngân hàng. Khung phương pháp tích hợp bốn thành phần: (1)~phân loại chủ đề ESG và (2)~phân loại mức độ hành động bằng mô hình PhoBERT kết hợp ràng buộc tri thức ký hiệu; (3)~liên kết bằng chứng; và (4)~chỉ số rủi ro ESG-Washing (EWRI).

Thực nghiệm trên 32.260 câu ESG từ báo cáo của 10 ngân hàng Việt Nam giai đoạn 2020--2024 cho thấy: (i)~mô hình phân loại đạt độ chính xác cao nhờ tích hợp ràng buộc tri thức; (ii)~phương pháp liên kết bằng chứng đề xuất vượt trội so với các phương pháp cơ sở dựa trên từ khóa; (iii)~chỉ số EWRI phân biệt có ý nghĩa giữa các ngân hàng; và (iv)~phân tích xu hướng cho thấy rủi ro ESG-washing giảm dần theo thời gian khi các công bố ngày càng nhấn mạnh các hành động cụ thể, có kiểm chứng. Kết quả khẳng định tính hiệu quả của phương pháp phân tích văn bản nội bộ trong phát hiện ESG-washing tại các thị trường mới nổi có hạn chế về dữ liệu xếp hạng bên ngoài.

\vspace*{1cm}
\textbf{Từ khóa}: ESG-washing, NLP, Neuro-Symbolic, Liên kết bằng chứng, Ngân hàng Việt Nam

\end{small}